\documentclass[../../../include/open-logic-section]{subfiles}

\begin{document}

\olfileid{sth}{spine}{foundation}

\olsection{بنیاد}

تنها \emph{تقریباً} کار را تمام کرده‌ایم---و هنوز \emph{کاملاً} به پایان
نرسیده‌ایم---زیرا هیچ چیز در~$\ZFminus$ تضمین نمی‌کند که \emph{هر}
مجموعه در یک~$V_\alpha$ باشد؛ یعنی هر مجموعه در مرحله‌ای تشکیل شده باشد.

اکنون، به یک معنای نسبتاً سرراستِ ریاضی، برایمان \emph{مهم نیست} که آیا
مجموعه‌هایی بیرون از سلسله‌مراتب وجود دارند یا نه. (اگر چنین مجموعه‌هایی
وجود داشته باشند، می‌توانیم صرفاً آن‌ها را نادیده بگیریم.) اما انگیزهٔ
\emph{مفهوم} مجموعه نزد ما این اندیشه بود که هر مجموعه در مرحله‌ای تشکیل
می‌شود (\stageshier{} را در~\olref[z][story]{sec} ببینید). پس می‌خواهیم
امکان وجود مجموعه‌هایی را که بیرون از سلسله‌مراتب قرار می‌گیرند منتفی
کنیم. در نتیجه باید اصل موضوع تازه‌ای بیفزاییم که تضمین کند هر مجموعه در
جایی از سلسله‌مراتب رخ می‌دهد.

چون~$V_\alpha$ها مراحل ما هستند، شاید به‌سادگی افزودن عبارت زیر را به‌عنوان
یک اصل موضوع در نظر بگیریم:

\begin{defish}
\emph{قاعده‌مندی.} $\forall A \exists \alpha\, A \subseteq V_\alpha$
\end{defish}

این رویکرد کاملاً معقولی می‌بود. بااین‌حال، به دلایلی که در بخش بعد
توضیح داده خواهند شد، به‌جای آن اصل موضوع بدیل زیر را اختیار می‌کنیم:

\begin{axiom}[بنیاد]
$(\forall A \neq \emptyset)(\exists B \in A)A \cap B = \emptyset$.
\end{axiom}

با قدری کوشش می‌توانیم در~$\ZFminus$ نشان دهیم بنیاد مستلزم قاعده‌مندی
است:
\begin{defn}
برای هر مجموعهٔ~$A$، بگذارید:
\begin{align*}
	\text{cl}_0(A) &= A,\\
	\text{cl}_{n+1}(A) &= \bigcup \text{cl}_n(A),\\
	\text{trcl}(A) &= \bigcup_{n < \omega} \text{cl}_{n}(A).
\end{align*}
$\text{trcl}(A)$ را \emph{بستار متعدیِ}~$A$ می‌نامیم.
\end{defn}
\noindent
نام «بستار متعدی» برازنده است:
\begin{prop}\ollabel{subsetoftrcl}
$A\subseteq\trcl{A}$ و~$\trcl{A}$ مجموعه‌ای متعدی است.
\end{prop}

\begin{proof}
آشکار است که~$A=\text{cl}_0(A)\subseteq\text{trcl}(A)$. همچنین اگر
$x\in b\in\trcl{A}$، آنگاه برای یک~$n$ داریم
$b\in\text{cl}_n(A)$، پس
$x\in\text{cl}_{n+1}(A)\subseteq\text{trcl}(A)$.
\end{proof}

\begin{lem}\ollabel{lem:TransitiveWellFounded}
اگر~$A$ مجموعه‌ای متعدی باشد، یک~$\alpha$ وجود دارد به‌طوری‌که
$A\subseteq V_\alpha$.
\end{lem}

\begin{proof}
با یادآوری تعریف «$\supstrict(X)$» از
\olref[ordinals][opps]{defsupstrict}، دو مجموعهٔ زیر را تعریف کنید:
\begin{align*}
	D  &= \Setabs{x \in A}{\forall \delta\ x \nsubseteq V_\delta}\\
	\alpha &= \supstrict\Setabs{\delta}{(\exists x \in A)
	(x \subseteq V_\delta \land (\forall \gamma \in \delta)x \nsubseteq V_\gamma)}
\end{align*}
فرض کنید~$D=\emptyset$. پس اگر~$x\in A$، یک~$\delta$ وجود دارد
به‌طوری‌که~$x\subseteq V_\delta$؛ و بنا بر خوش‌ترتیبی اعداد ترتیبی،
$(\forall \gamma\in\delta)x\nsubseteq V_\gamma$. در نتیجه
$\delta\in\alpha$ و بنا بر
\olref[spine][Valphabasic]{Valphabasicprops}، $x\in V_\alpha$. پس
$A\subseteq V_\alpha$، چنان‌که می‌خواستیم.

بنابراین کافی است نشان دهیم~$D=\emptyset$. برای برهان خلف خلافش را فرض
کنید. بنا بر بنیاد، یک~$B\in D\subseteq A$ وجود دارد به‌طوری‌که
$D\cap B=\emptyset$. اگر~$x\in B$، چون~$A$ متعدی است، $x\in A$؛ و چون
$x\notin D$، نتیجه می‌شود~$\exists\delta\ x\subseteq V_\delta$. اکنون
بگذارید
\[
	\beta = \supstrict\Setabs{\delta}{(\exists x \in B)(x \subseteq V_\delta \land (\forall \gamma < \delta)x \nsubseteq V_\gamma)}.
\]
همانند قبل، $B\subseteq V_\beta$؛ و این با~$B\in D$ تناقض دارد.
\end{proof}

\begin{thm}\ollabel{zfentailsregularity}
قاعده‌مندی برقرار است.
\end{thm}

\begin{proof}
$A$ را ثابت بگیرید. بنا بر~\olref{subsetoftrcl}،
$A\subseteq\trcl{A}$ و~$\trcl{A}$ متعدی است. پس بنا بر
\olref{lem:TransitiveWellFounded} یک~$\alpha$ وجود دارد به‌طوری‌که
$A\subseteq\trcl{A}\subseteq V_\alpha$.
\end{proof}

این نتایج نشان می‌دهند~$\ZFminus$ شرطیِ
$\emph{بنیاد}\Rightarrow\emph{قاعده‌مندی}$ را اثبات می‌کند. در
\olref[spine][rank]{zfminusregularityfoundation} نشان خواهیم داد
$\ZFminus$ شرطیِ
$\emph{قاعده‌مندی}\Rightarrow\emph{بنیاد}$ را نیز اثبات می‌کند.
بنابراین بنیاد و قاعده‌مندی (نسبت به~$\ZFminus$) \emph{هم‌ارزند}. اما این
بدین معناست که با فرض~$\ZFminus$ می‌توانیم بنیاد را با هم‌ارزی‌اش با
قاعده‌مندی توجیه کنیم؛ و قاعده‌مندی را نیز بی‌درنگ بر پایهٔ
\stageshier{} توجیه می‌کنیم.

\end{document}
